\documentclass[10pt,a4paper]{amsart}
\usepackage{amsmath,amssymb,amsthm,graphicx}
\usepackage{verbatim}
\newenvironment{reference}{\comment}{\endcomment}
%\newenvironment{reference}{}{}
\newenvironment{slogan}{\comment}{\endcomment}
\newenvironment{history}{\comment}{\endcomment}

% For commutative diagrams we use Xy-pic
\usepackage[all]{xy}

% We use 2cell for 2-commutative diagrams.
\xyoption{2cell}
% Theorem environments.
%
\theoremstyle{plain}
\newtheorem{theorem}[subsection]{Theorem}
\newtheorem{proposition}[subsection]{Proposition}
\newtheorem{lemma}[subsection]{Lemma}

\theoremstyle{definition}
\newtheorem{definition}[subsection]{Definition}
\newtheorem{example}[subsection]{Example}
\newtheorem{exercise}[subsection]{Exercise}
\newtheorem{situation}[subsection]{Situation}

\theoremstyle{remark}
\newtheorem{remark}[subsection]{Remark}
\newtheorem{remarks}[subsection]{Remarks}

\numberwithin{equation}{subsection}

% Macros
%
\def\lim{\mathop{\mathrm{lim}}\nolimits}
\def\colim{\mathop{\mathrm{colim}}\nolimits}
\def\Spec{\mathop{\mathrm{Spec}}}
\def\Hom{\mathop{\mathrm{Hom}}\nolimits}
\def\Ext{\mathop{\mathrm{Ext}}\nolimits}
\def\SheafHom{\mathop{\mathcal{H}\!\mathit{om}}\nolimits}
\def\SheafExt{\mathop{\mathcal{E}\!\mathit{xt}}\nolimits}
\def\Sch{\mathit{Sch}}
\def\Mor{\mathop{\mathrm{Mor}}\nolimits}
\def\Ob{\mathop{\mathrm{Ob}}\nolimits}
\def\Sh{\mathop{\mathit{Sh}}\nolimits}
\def\NL{\mathop{N\!L}\nolimits}
\def\CH{\mathop{\mathrm{CH}}\nolimits}
\def\proetale{{pro\text{-}\acute{e}tale}}
\def\etale{{\acute{e}tale}}
\def\QCoh{\mathit{QCoh}}
\def\Ker{\mathop{\mathrm{Ker}}}
\def\Im{\mathop{\mathrm{Im}}}
\def\Coker{\mathop{\mathrm{Coker}}}
\def\Coim{\mathop{\mathrm{Coim}}}

% Boxtimes
%
\DeclareMathSymbol{\boxtimes}{\mathbin}{AMSa}{"02}

%
% Macros for moduli stacks/spaces
%
\def\QCohstack{\mathcal{QC}\!\mathit{oh}}
\def\Cohstack{\mathcal{C}\!\mathit{oh}}
\def\Spacesstack{\mathcal{S}\!\mathit{paces}}
\def\Quotfunctor{\mathrm{Quot}}
\def\Hilbfunctor{\mathrm{Hilb}}
\def\Curvesstack{\mathcal{C}\!\mathit{urves}}
\def\Polarizedstack{\mathcal{P}\!\mathit{olarized}}
\def\Complexesstack{\mathcal{C}\!\mathit{omplexes}}
% \Pic is the operator that assigns to X its picard group, usage \Pic(X)
% \Picardstack_{X/B} denotes the Picard stack of X over B
% \Picardfunctor_{X/B} denotes the Picard functor of X over B
\def\Pic{\mathop{\mathrm{Pic}}\nolimits}
\def\Picardstack{\mathcal{P}\!\mathit{ic}}
\def\Picardfunctor{\mathrm{Pic}}
\def\Deformationcategory{\mathcal{D}\!\mathit{ef}}
\usepackage{cleveref}
\crefname{originalreference}{original reference}{original references}
\Crefname{originalreference}{Original reference}{Original references}
\makeatletter
\newlabel{definition-formal-modification}{{0GDK}{source}{}{sourcepointer-1}{}}
\newlabel{definition-formal-modification@cref}{{[originalreference][0][]0GDK}{[1][0][]source}{}{}{}}
\makeatother
\crefname{originaldefinition}{original definition}{original definitions}
\Crefname{originaldefinition}{Original definition}{Original definitions}
\makeatletter
\newlabel{formal-spaces-definition-adic-morphism}{{0AQ3}{source}{}{sourcepointer-2}{}}
\newlabel{formal-spaces-definition-adic-morphism@cref}{{[originaldefinition][0][]0AQ3}{[1][0][]source}{}{}{}}
\makeatother
\crefname{originallemma}{original lemma}{original lemmas}
\Crefname{originallemma}{Original lemma}{Original lemmas}
\makeatletter
\newlabel{formal-spaces-lemma-map-completions-representable}{{0APV}{source}{}{sourcepointer-3}{}}
\newlabel{formal-spaces-lemma-map-completions-representable@cref}{{[originallemma][0][]0APV}{[1][0][]source}{}{}{}}
\makeatother
\crefname{originallemma}{original lemma}{original lemmas}
\Crefname{originallemma}{Original lemma}{Original lemmas}
\makeatletter
\newlabel{formal-spaces-lemma-map-into-algebraic-space}{{0AQH}{source}{}{sourcepointer-4}{}}
\newlabel{formal-spaces-lemma-map-into-algebraic-space@cref}{{[originallemma][0][]0AQH}{[1][0][]source}{}{}{}}
\makeatother
\crefname{originalreference}{original reference}{original references}
\Crefname{originalreference}{Original reference}{Original references}
\makeatletter
\newlabel{lemma-modification-gives-formal-modification}{{0GDM}{source}{}{sourcepointer-5}{}}
\newlabel{lemma-modification-gives-formal-modification@cref}{{[originalreference][0][]0GDM}{[1][0][]source}{}{}{}}
\makeatother
\crefname{originalreference}{original reference}{original references}
\Crefname{originalreference}{Original reference}{Original references}
\makeatletter
\newlabel{lemma-output-quasi-compact}{{0ARU}{source}{}{sourcepointer-6}{}}
\newlabel{lemma-output-quasi-compact@cref}{{[originalreference][0][]0ARU}{[1][0][]source}{}{}{}}
\makeatother
\crefname{originalreference}{original reference}{original references}
\Crefname{originalreference}{Original reference}{Original references}
\makeatletter
\newlabel{lemma-output-separated}{{0ARW}{source}{}{sourcepointer-7}{}}
\newlabel{lemma-output-separated@cref}{{[originalreference][0][]0ARW}{[1][0][]source}{}{}{}}
\makeatother
\crefname{originallemma}{original lemma}{original lemmas}
\Crefname{originallemma}{Original lemma}{Original lemmas}
\makeatletter
\newlabel{spaces-cohomology-lemma-check-proper-dvr}{{0ARK}{source}{}{sourcepointer-8}{}}
\newlabel{spaces-cohomology-lemma-check-proper-dvr@cref}{{[originallemma][0][]0ARK}{[1][0][]source}{}{}{}}
\makeatother
\crefname{originalremark}{original remark}{original remarks}
\Crefname{originalremark}{Original remark}{Original remarks}
\makeatletter
\newlabel{spaces-cohomology-remark-variant}{{0ARL}{source}{}{sourcepointer-9}{}}
\newlabel{spaces-cohomology-remark-variant@cref}{{[originalremark][0][]0ARL}{[1][0][]source}{}{}{}}
\makeatother
\crefname{originalreference}{original reference}{original references}
\Crefname{originalreference}{Original reference}{Original references}
\makeatletter
\newlabel{theorem-dilatations-general}{{0ARB}{source}{}{sourcepointer-10}{}}
\newlabel{theorem-dilatations-general@cref}{{[originalreference][0][]0ARB}{[1][0][]source}{}{}{}}
\makeatother
\title{The algebraized map is proper exactly when the formal map is a formal modification}
\author{The Stacks Project Authors}
\date{stacks-project:restricted.tex:local-snapshot-20261012; GNU FDL 1.2 or later}
\begin{document}
\maketitle
\noindent\textit{Editorial note.} All delivered mathematical text is the author's own, copied from the pinned Stacks Project chapter file restricted.tex; no statement or proof step was rewritten, reordered or removed.\par\medskip
\noindent\textit{Editorial note.} The counted proof is the author's complete argument: the proof is complete: both directions are given and the three DVR cases are each carried out\par\medskip
\noindent\textit{Editorial note.} The label lemma-output-proper names the counted statement, and the proof environment immediately following it in the source is delivered verbatim as its proof.\par\medskip
\noindent\textit{Editorial note.} The counted statement is carried by the author's own lemma environment in the chapter source.\par\medskip
\noindent\textit{Allowed standard prerequisites (editorial).}\begin{itemize}
\item the dilatation correspondence of the chapter between maps f : X' to X and rig-etale maps g : W to X\_\{/T\}
\item formal modifications
\item the DVR lifting criterion for properness of morphisms of algebraic spaces and its rig-surjective variant
\item the quasi-compactness and separatedness output lemmas for the correspondence
\item algebraization of adic morphisms of formal spectra into algebraic spaces
\end{itemize}\medskip
\ifdefined\hypertarget\hypertarget{sourcepointer-1}{}\fi
\noindent\textit{Source reference.} Original reference 0GDK: sources/stacks-restricted/pinned-original-source.tex, source lines 6275--6275. Prerequisite pointer invoked at proof line(s) 7283; it names a standard fact of the same chapter that the counted argument uses, and every step of the counted proof built on it is retained verbatim in the delivered body.\par\medskip
\ifdefined\hypertarget\hypertarget{sourcepointer-2}{}\fi
\noindent\textit{Source reference.} Original definition 0AQ3: sources/stacks-chow/tags.txt, source lines 13496--13496. Prerequisite pointer invoked at proof line(s) 7326; it names a standard fact of a sibling chapter that the counted argument uses, and every step of the counted proof built on it is retained verbatim in the delivered body.\par\medskip
\ifdefined\hypertarget\hypertarget{sourcepointer-3}{}\fi
\noindent\textit{Source reference.} Original lemma 0APV: sources/stacks-chow/tags.txt, source lines 13488--13488. Prerequisite pointer invoked at proof line(s) 7325; it names a standard fact of a sibling chapter that the counted argument uses, and every step of the counted proof built on it is retained verbatim in the delivered body.\par\medskip
\ifdefined\hypertarget\hypertarget{sourcepointer-4}{}\fi
\noindent\textit{Source reference.} Original lemma 0AQH: sources/stacks-chow/tags.txt, source lines 13510--13510. Prerequisite pointer invoked at proof line(s) 7332; it names a standard fact of a sibling chapter that the counted argument uses, and every step of the counted proof built on it is retained verbatim in the delivered body.\par\medskip
\ifdefined\hypertarget\hypertarget{sourcepointer-5}{}\fi
\noindent\textit{Source reference.} Original reference 0GDM: sources/stacks-restricted/pinned-original-source.tex, source lines 6346--6346. Prerequisite pointer invoked at proof line(s) 7288; it names a standard fact of the same chapter that the counted argument uses, and every step of the counted proof built on it is retained verbatim in the delivered body.\par\medskip
\ifdefined\hypertarget\hypertarget{sourcepointer-6}{}\fi
\noindent\textit{Source reference.} Original reference 0ARU: sources/stacks-restricted/pinned-original-source.tex, source lines 7117--7117. Prerequisite pointer invoked at proof line(s) 7290; it names a standard fact of the same chapter that the counted argument uses, and every step of the counted proof built on it is retained verbatim in the delivered body.\par\medskip
\ifdefined\hypertarget\hypertarget{sourcepointer-7}{}\fi
\noindent\textit{Source reference.} Original reference 0ARW: sources/stacks-restricted/pinned-original-source.tex, source lines 7198--7198. Prerequisite pointer invoked at proof line(s) 7290; it names a standard fact of the same chapter that the counted argument uses, and every step of the counted proof built on it is retained verbatim in the delivered body.\par\medskip
\ifdefined\hypertarget\hypertarget{sourcepointer-8}{}\fi
\noindent\textit{Source reference.} Original lemma 0ARK: sources/stacks-chow/tags.txt, source lines 13548--13548. Prerequisite pointer invoked at proof line(s) 7294; it names a standard fact of a sibling chapter that the counted argument uses, and every step of the counted proof built on it is retained verbatim in the delivered body.\par\medskip
\ifdefined\hypertarget\hypertarget{sourcepointer-9}{}\fi
\noindent\textit{Source reference.} Original remark 0ARL: sources/stacks-chow/tags.txt, source lines 13549--13549. Prerequisite pointer invoked at proof line(s) 7295; it names a standard fact of a sibling chapter that the counted argument uses, and every step of the counted proof built on it is retained verbatim in the delivered body.\par\medskip
\ifdefined\hypertarget\hypertarget{sourcepointer-10}{}\fi
\noindent\textit{Source reference.} Original reference 0ARB: sources/stacks-restricted/pinned-original-source.tex, source lines 7015--7015. Prerequisite pointer invoked at proof line(s) 7280; it names a standard fact of the same chapter that the counted argument uses, and every step of the counted proof built on it is retained verbatim in the delivered body.\par\medskip
\begin{lemma}
\label{lemma-output-proper}
With assumptions and notation as in Theorem \ref{theorem-dilatations-general}
let $f : X' \to X$ correspond to $g : W \to X_{/T}$.
Then $f$ is proper if and only if $g$ is a formal modification
(Definition \ref{definition-formal-modification}).
\end{lemma}

\begin{proof}
If $f$ is proper, then $g$ is a formal modification by
Lemma \ref{lemma-modification-gives-formal-modification}.
Assume $g$ is a formal modification. By
Lemmas \ref{lemma-output-quasi-compact} and \ref{lemma-output-separated}
we see that $f$ is quasi-compact and separated.

\medskip\noindent
By Cohomology of Spaces, Lemma \ref{spaces-cohomology-lemma-check-proper-dvr}
and Remark \ref{spaces-cohomology-remark-variant}
it suffices to show that given any commutative diagram
$$
\xymatrix{
\Spec(K) \ar[r] \ar[d] & X' \ar[d]^f \\
\Spec(R) \ar[r]^p \ar@{-->}[ru] & X
}
$$
where $R$ is a complete discrete valuation ring with fraction field $K$,
there is a dotted arrow making the diagram commute.
There are three cases:
Case I: $p(\Spec(R)) \subset U$. This case is trivial
because $U' \to U$ is an isomorphism.
Case II: $p$ maps $\Spec(R)$ into $T$. This case follows
from our assumption that $g : W \to X_{/T}$ is proper. Namely,
if $Z$ denotes the reduced induced closed subspace structure
on $T$, then $p$ factors through $Z$ and
$$
W \times_{X_{/T}} Z = X' \times_X Z \longrightarrow Z
$$
is proper by assumption which implies we get the lifting property by
Cohomology of Spaces, Lemma \ref{spaces-cohomology-lemma-check-proper-dvr}
applied to the displayed arrow. Case III: $p(\Spec(K))$ is not in $T$
but $p$ maps the closed point of $\Spec(R)$ into $T$. In this case
the corresponding morphism
$$
p_{/T} : \text{Spf}(R) \longrightarrow X'_{/T} = W
$$
is an adic morphism (by
Formal Spaces, Lemma
\ref{formal-spaces-lemma-map-completions-representable} and
Definition \ref{formal-spaces-definition-adic-morphism}).
Hence our assumption that $g : W \to X_{/T}$ be rig-surjective
implies we can lift $g_{/T}$ to a morphism
$\text{Spf}(R') \to W = X'_{/T}$
for some extension of complete discrete valuation rings $R \subset R'$.
Algebraizing the composition $\text{Spf}(R') \to X'$ using
Formal Spaces, Lemma \ref{formal-spaces-lemma-map-into-algebraic-space}
we find a morphism $\Spec(R') \to X'$ lifting $p$ as desired.
\end{proof}

\end{document}
